\chapter{绪论}

%章节内容可以根据需要增添。

\section{研究背景及意义}
合成孔径雷达（Synthetic Aperture Radar，SAR）是一种基于电磁波主动探测的微波成像技术，通过向观测目标发射微波信号并接收后向散射回波，
结合平台运动构建虚拟合成孔径，突破天线孔径限制实现高分辨率成像，
其独特的微波传输特性使其能够穿透云层，在昼夜交替、恶劣天气等光学遥感无法工作的场景下保持稳定观测能力，
自 20 世纪 50 年代问世以来，已成为军民领域重要的对地观测手段
\cite{ref-1-1,李春升2019高分辨率星载,el2013target}。

随着SAR成像技术的发展，其应用场景从早期军事侦察逐步演进至战场态势感知、制导毁伤等核心领域，并成功拓展至民用范畴。
相较于光学成像手段，SAR能够在夜间、云雾遮挡以及复杂气象条件下稳定获取目标及其周边场景信息，
因而在装甲车辆、舰船、飞机等战术目标的检测、跟踪与识别中具有重要应用价值。
随着无人机、导弹、卫星等平台的快速发展，前端感知系统对目标信息获取的实时性、自动化和鲁棒性提出了更高要求，
SAR图像目标识别逐渐成为提升战场态势感知能力和支撑自主决策的重要技术环节。
受 SAR 成像机理的固有特性影响，其图像普遍存在大量相干斑噪声，且同类目标在不同俯仰角、
方位角和成像条件下的特征差异显著。传统人工判读模式不仅需要投入大量人力物力成本，
而且难以满足复杂战场环境下快速响应和无人平台自主应用的现实需求，
因此旨在减少乃至替代人工解译的SAR图像自动目标识别（Automatic Target Recognition，ATR）
技术得以快速发展并获得广泛应用\cite{el2013target,kechagias2021automatic}。

传统的SAR-ATR主要由3个阶段组成：检测、鉴别和分类，随着SAR成像质量的提升和信息的丰富,解译任务开始朝着更细粒度的方向发展\cite{罗汝2023sar,何奇山2024面向}。
其中SAR目标的细粒度分类是最为核心的一个问题。近年来，
随着人工智能（Artificial Intelligence，AI）技术，尤其是深度神经网络（Deep Neural Network，DNN）的发展，
基于数据驱动和神经网络模型的SAR目标识别取得了卓越的性能\cite{chen2016target,王俊2018深度学习在雷达中的研究综述}。
现有的目标识别模型在经典的SAR 目标识别数据集 MSTAR（Moving and Stationary Target Acquisition and Recognition）
的标准工作条件（Standard Operating Condition，SOC）上已能够达到99\%的识别准确率\cite{li2023hierarchical}。
然而，由于合成孔径雷达与深度学习技术的固有特性，数据驱动的SAR-ATR的实际应用仍面临以下挑战：

\textbf{(1)实测训练数据获取瓶颈}：深度学习模型的训练依赖大规模高质量实测 SAR 图像数据集。
然而，实测 SAR 图像的获取需调配高精度专用成像设备，标注过程需依托 SAR 图像处理领域的专业知识，
导致数据获取与标注的经济成本与时间成本居高不下。此外，SAR 图像的成像质量对平台参数、目标朝向及观测角度等因素具有高度敏感性：
若成像平台参数或作业环境条件发生变化，为保障模型的识别性能达标，需重新开展训练数据采集工作，
进一步推高了数据获取的综合成本\cite{yang2024transfer}。更为关键的是，在部分特殊应用场景（如军事侦察、涉密目标监测等）中，
目标往往具备非合作性、低暴露性及强隐蔽性特点，导致在真实应用场景下获取足量的高质量训练数据面临严峻挑战\cite{何奇山2024面向}。

\textbf{(2)数据分布过拟合风险突出}：深度学习方法通常基于特定概率分布的训练数据学习目标表征，
并默认训练数据与测试数据满足独立同分布（Independent and Identically Distributed，i.i.d）假设。
当模型用于训练的源域（Source Domain）数据和待测的目标域（Target Domain）数据之间分布差异较大时，
模型学习到的表征也将出现显著差异,从而导致测试性能大幅下降。
该现象被称为 “域偏移”（Domain Shift）。在此情况下，模型易对训练数据中由特定 SAR 成像参数（如波段、极化方式、俯仰角等）
诱导的专属冗余特征发生过拟合，进而导致模型在非同源测试数据上的泛化性能与识别精度显著退化\cite{yang2024transfer}。

\textbf{(3)边缘设备部署面临算力约束}：SAR 图像目标识别的实际应用场景（如战场侦察、精确制导、低空监测等）常要求模型部署于无人机、
导弹、便携式探测终端等边缘计算节点。此类设备受限于体积、功耗及制造成本约束，硬件算力资源
存在先天不足，难以承载深度学习模型的高复杂度计算需求。一方面，模型推理过程中大量的矩阵运算、特征映射及非线性变换操作会导致边缘设备计算延迟过高，
无法满足实时识别的应用需求；另一方面，高算力消耗会急剧缩短边缘设备的续航时间，限制其在无外接电源场景下的持续工作能力。
此外，复杂模型通常伴随较大的参数量与存储占用，而边缘设备的存储空间往往有限，进一步加剧了模型部署的难度\cite{xu2022lite}。
上述算力与存储约束导致许多在服务器端表现优异的 SAR 图像识别模型难以落地应用，严重限制了相关技术的实际推广价值。



随着计算机辅助设计（CAD）和电磁仿真等相关技术的发展，通过仿真手段生成大量高质量SAR图像已成为可能。
仿真手段能灵活、高效地修改SAR成像平台运动条件和目标姿态，并且可以方便地获取目标详细的标签信息，
因此利用仿真SAR数据替代实测数据训练模型正逐渐成为一种可行且低成本的方案。典型的SAR图像仿真流程包括构建目标的三维模型，
并将电磁计算与计算机图形学方法相结合，以生成仿真 SAR 图像。对于某些特定目标，仿真图像能够实现相对逼真的视觉效果\cite{ref-2015-3-Dong,ref-2021-12-Hu}。
然而，以往的研究指出，仿真 SAR 图像与实测图像在背景纹理、目标结构等方面不可避免地存在差异\cite{lewis2019sar,inkawhich2021bridging}。这些差异的产生，
一方面源于仿真算法本身的不足，另一方面则是由于真实场景中存在复杂的环境噪声。这表明仿真图像与真实图像之间仍存在显著的域偏移现象，
如何解决从仿真到实测的域偏移问题成为了利用仿真SAR数据研究的重点。
有鉴于此，本文拟重点针对面向仿真到实测迁移的跨域SAR目标识别方法展开研究，
旨在充分探究在现实环境中实测数据保障受限的情况下如何利用仿真数据训练并迁移SAR目标识别模型。

此外，在国产AI芯片应用逐步推进的背景下，相关平台在自主可控、供应链稳定性和敏感场景部署等方面受到越来越多关注\cite{钱贵明2025ai}。
目前科研中常用于深度学习部署加速的平台，例如NVIDIA公司的Jetson TX系列，其底层推理框架和开发生态主要依赖国外软硬件平台。
对于信息安全、供应链连续性或长期可维护性要求较高的应用场景，这种依赖会带来一定限制。
近年来，华为、瑞芯微、寒武纪等国内厂商持续推进AI芯片研发，但在生态成熟度、工具链完善程度和市场规模等方面仍与成熟进口平台存在差距。
因此，围绕国产AI芯片开展稳定可用的部署技术研究，对SAR目标识别等应用场景具有现实意义。

在本文研究的 SAR 图像目标识别问题中，“域”可理解为由不同传感器参数、观测角度、成像条件或数据获取平台形成的、统计分布存在显著差异的 SAR 数据集合。
在实际应用中，训练数据与测试数据往往难以保持同源：一方面，实测 SAR 数据获取成本高、周期长，且受任务时机和场景条件限制明显；另一方面，平台更替、任务转换和环境变化又会持续引入新的数据分布。
因此，跨域问题并非 SAR 目标识别中的偶发现象，而是制约模型走向实用化的核心瓶颈之一。
在本文后续章节中，源域主要指电磁仿真生成的SAR数据，目标域主要指实际采集的实测SAR数据，本文所关注的跨域迁移主要对应“仿真到实测”的迁移过程。
围绕仿真到实测迁移这一核心任务，本文进一步将跨域 SAR 目标识别中的典型困难归纳为“少量实测数据可见（目标域数据少样本）”和“实测数据不可见”两类场景。

本文所定义的“目标域数据少样本”是指：在模型训练阶段，仅能获取少量与待测实测数据分布一致的无标签 SAR 图像。
这一场景对应于实际任务中虽然能够事先接触部分实测数据，但其数量较为有限、且难以完成人工标注的情况。
在此条件下，少量目标域样本虽然能够提供一定的分布参考，却不足以支撑传统深度学习方法完成稳定训练，
这就使得如何利用有限实测信息缓解仿真与实测之间的分布偏移，成为提升模型实际可用性的重要问题。
因此，有必要研究更适用于少样本目标域条件的（Domain Adaptation, DA）方法，在尽可能少依赖实测数据的前提下提升跨域识别性能。

本文所定义的“目标域数据不可见”是指：在模型训练阶段，完全无法获取与待测实测数据分布一致的 SAR 图像。
这一场景更接近实际部署中的极端条件，例如面对未知平台、陌生区域或临机任务时，模型往往无法提前获得任何目标域样本。
此时，模型既不能直接感知实测数据分布，也无法借助实测数据进行对齐或校正，只能依赖仿真数据中更具普适性的判别信息完成识别。
因此，如何从源域中挖掘稳定的域不变特征、提升模型对未知分布的泛化能力，成为保障跨域 SAR 目标识别鲁棒性的关键问题，
这也使得域泛化(Domain Generalization, DG)研究具有突出的理论意义与实用价值。

除了识别性能之外，跨域 SAR 目标识别方法还必须面向实际部署需求。
考虑到相关模型最终往往需要运行于无人平台或边缘设备，算力、功耗和存储资源均较为受限，
若缺乏轻量化处理与国产平台适配，即使模型具有较好的跨域识别能力，也难以真正落地应用。
因此，在研究跨域识别方法的同时，进一步开展模型压缩、推理优化与国产 AI 芯片部署验证，同样具有重要现实意义。

综上，本文面向仿真到实测迁移场景下的 SAR 目标识别需求，
围绕“少量实测数据可见”和“实测数据不可见”两类典型场景，分别研究适用于有限实测信息条件的域适应方法和面向未知实测域的域泛化方法，
重点解决利用电磁仿真 SAR 图像训练模型时面临的仿真与实测之间的域偏移问题；
在此基础上，开展模型轻量化与国产 AI 芯片平台部署验证，
构建从跨域识别方法研究到端侧应用验证的研究链路，为复杂战场环境下的 SAR 目标识别任务提供理论依据与技术支撑。

\section{国内外研究现状}
本节主要从跨域迁移学习技术的研究现状、跨域SAR目标识别的研究现状两个层面展开。
\subsection{跨域迁移学习技术的研究现状}
广义而言，迁移学习（Transfer Learning）旨在利用已有任务（源域）中习得的知识、模型参数或特征结构，
辅助并优化新任务（目标域）的学习过程。作为一种典型的数据驱动范式，深度学习高度依赖于训练集与测试集满足独立同分布的假设。
然而在实际应用中，由于训练数据与测试数据往往存在概率分布差异（即域偏移），
导致模型极易过拟合于源域的数据分布，进而在其他分布的数据上出现严重的性能退化。
因此，如何有效克服域偏移、实现源域知识向目标域的鲁棒迁移，从而提升模型在目标数据上的泛化能力，
构成了跨域迁移学习的核心研究目标\cite{zhuang2020comprehensive}。

域适应（Domain Adaptation, DA）作为跨域迁移学习的重要分支，旨在联合利用带标签的源域数据与无标签（或仅含少量标签）的目标域数据进行模型训练，
以最大化给定目标域上的测试性能。其核心机制通常是通过对齐源域与目标域的特征分布来缩小域间偏移，
进而有效改善目标域标签匮乏与域间分布不匹配等问题\cite{wang2018deep}。

尽管域适应在一定程度上缓解了分布偏移问题，但其前提是要求在训练阶段必须能够访问目标域数据，
这在许多实际场景中带来了较大限制。
为突破这一瓶颈，研究人员进一步提出了域泛化（Domain Generalization, DG）。
域泛化的目标是从单一或多个分布不同但内在相关的源域中学习到领域不变特征（Domain-invariant Features），
从而使模型在完全不可见（Unseen）的全新目标域上仍具备较好的泛化性能\cite{zhou2022domain}。

结合不同应用场景的实际需求，在SAR目标识别领域引入域适应与域泛化等跨域迁移学习技术，
应对训练与测试数据非同源引发的跨域识别难题，具有重要的理论和工程意义。
域适应与域泛化的实现机制对比如图\ref{fig:DA_DG}所示，本节后续将分别对域适应与域泛化技术的研究现状进行详细梳理。






\begin{figure}[htbp]  
  \centering          
  \subfloat[域适应实现]{\label{域适应实现} \includegraphics[width=0.44\textwidth]{chap01/域适应实现.png}} \hspace{0.6em}
  \subfloat[域泛化实现]{\label{域泛化实现} \includegraphics[width=0.48\textwidth]{chap01/域泛化实现.png}} 
  \caption{域适应与域泛化工作机制的对比示意图} 
  \label{fig:DA_DG}                  
\end{figure}

\subsubsection{域适应的研究现状}

域适应任务可定义如下：源域数据集$D_s = \{(x_i^s, y_i^s)\}_{i=1}^{n_s}$由$n_s$个源域样本$x_i^s$及其对应的标签$y_i^s$构成，
其中源域训练数据均带有标注信息；目标域数据集$D_t = \{x_i^t\}_{i=1}^{n_t}$由$n_t$个目标域样本$x_i^t$构成，
且目标域训练数据仅有少量或无标注信息。域适应的核心任务是学习一个映射函数$f: \mathcal{X} \to \mathcal{Y}$，
借助带有标注的源域数据集$D_s$和无标注的目标域数据集$D_t$，完成模型训练，以实现对目标域数据集$D_t$中样本的标签预测$\hat{y}_i^t = f(x_i^t)$，
需注意的是，源域数据$D_s$与目标域数据$D_t$的概率特征分布存在固有差异
（即$P_s(X) \neq P_t(X)$，其中$P_s(X)$和$P_t(X)$分别表示源域和目标域的特征分布）。
基于深度学习的域适应方法主要可划分为四类：基于差异的方法、基于对抗的方法、基于重构的方法以及基于生成的方法。


（1）基于差异的方法

基于差异的域适应方法主要指的是通过设计特征分布距离度量来衡量源域与目标域之间的分布差异，并通过最小化该距离度量来减小域偏移，
从而实现域适应。典型代表如Deep CORAL\cite{deepcoral}在深层网络中加入相关对齐损失，
对齐源域与目标域特征的协方差矩阵，从而在不改变网络结构的前提下实现二阶统计层面的域对齐。
Long等人\cite{RTN}提出的残差迁移网络（Residual Transfer Network，RTN）则在多层特征上引入基于极大均值差异（Maximum Mean Discrepancy，MMD）的核均值匹配，
同时将目标分类器建模为源分类器的残差修正，联合完成特征与分类器的迁移。
这类方法结构简单、训练稳定，易于与任意卷积主干结合，因此最早被成体系地应用于图像分类任务。
随后大量工作在其基础上进行扩展：一类是将对齐操作直接嵌入网络归一化层，
例如进行特征白化（Feature‑Whitening）\cite{roy2019unsupervised}或设计域对齐层（Domain Alignment Layers）\cite{multidial}，
通过在批归一化层位置实现分布标准化来隐式减小域差异；
另一类则面向视觉任务的细粒度语义结构，引入类条件对齐、原型网络或对比学习\cite{cdcl,clusteralignment,orthogonalclustered}，
在特征空间按类别对齐条件分布，缓解早期全局对齐容易造成“错误类对齐”和决策边界模糊的问题，
总体来看，基于差异的方法已经从“全局统计对齐”发展为结合矩匹配、归一化、类条件原型与对比损失的系统家族，并与对抗、自监督训练等机制形成互补。


（2）基于对抗的方法

基于对抗的域适应框架以域判别器与特征提取器之间的对抗博弈为核心：域判别器通过最小化域分类损失学习区分输入来自源域还是目标域，
而特征提取器则通过最大化该损失（通常借助梯度反转层实现）逼迫中间特征“域不可辨”，
在同时最小化源域任务损失的约束下，得到既域不变又类别可分的表示。
Ganin和Lempitsky提出\cite{ganin2015unsupervised}的域对抗网络（Domain Adversarial Neural Network，DANN），首次将对抗学习思想与深度域适应方法相结合，
给出了一个与主干结构基本正交、可直接插入任意骨干网络的通用“域对齐模块”。后续的许多工作基本都沿着“对抗对齐 + 任务网络”的方向开展拓展。
Tzeng等人\cite{tzeng2017adversarial}提出的对抗判别域适应（Adversarial Discriminative Domain Adaptation，ADDA）
方法，突破了DANN 特征提取器与判别器共享参数的限制，采用非共享参数架构，分阶段训练以强化目标域特征的特异性学习，
使得模型能够学习更多领域特有特征，从而提升复杂域偏移下的适配能力。
Long等人\cite{long2018conditional}则借鉴条件生成对抗网络（Conditional Generative Adversarial Networks，CGAN）\cite{mirza2014conditional}
的对抗学习框架，构建了条件域对抗网络（Conditional Domain Adversarial Networks，CDAN），
其核心设计思路在于实现源域与目标域条件分布的精准对齐。与 DANN 直接将特征提取器输出的特征向量输入域判别器的做法不同，
CDAN 将输入形式调整为任务分类器输出的类别预测向量与特征提取器提取的特征向量的外积形式，
这一改进使得模型能够融入目标域样本的类别细粒度信息，最终进一步增强了模型的识别性能。
总体上，“对抗对齐 + 任务网络”这种模式已成为域适应领域的常见范式\cite{pei2018multi,hoffman2018cycada,tang2020discriminative,cao2018partial}，
并在图像分类、目标检测、语义分割等多个视觉任务中得到广泛应用\cite{hassanpour2023survey}，但仍存在对目标域数据需求较高、训练过程不稳定和易出现模式崩溃等问题。



（3）基于重构的方法

基于重构的域适应方法核心思想在于，通过自编码器编码获得具有一定迁移性的特征表达，借助“由特征重构输入”的约束机制，
迫使模型学习跨域一致的结构化特征（即域不变特征），进而间接缩小源域与目标域的分布差异。
Ghifary等人\cite{ghifary2016deep}首次提出深度重构分类网络（Deep Reconstruction Classification Networks，DRCN），
该网络设计共享编码器以提取源域与目标域的特征，利用源域标注样本训练任务分类器以保障特征的判别能力，
同时采用目标域图像训练目标域解码器完成重构任务。通过目标域特征到原始输入的重构过程，
编码器输出的特征需同时满足“源域分类任务需求”与“目标域结构保留需求”，最终生成跨域通用特征，间接实现域分布对齐。
Bousmalis等人\cite{bousmalis2016domain}指出，模型提取的特征可拆解为跨域共有特征与域专属特征，为在域适应过程中剔除冗余的域专属特征，
提出域分离网络（Domain Separation Network，DSN）。该网络通过目标域专属编码器、源域专属编码器及共享编码器，
分别提取目标域专属特征、源域专属特征与跨域共有特征，再利用共享解码器融合专属特征与共有特征完成图像重构，实现共有特征与专属特征的有效分离。
后续研究中，部分基于重构的域适应方法进一步与对抗学习机制相结合。
例如，Peng等人\cite{peng2019domain}提出了一种深度对抗解耦自编码器（Deep Adversarial Disentangled Autoencoder，DADA），
针对多目标域适应任务，通过对抗训练将特征解纠缠为“域相关成分”与“任务相关成分”；
通过域判别器仅对“域相关成分”进行优化，实现仅对齐任务相关成分的精准域适应，最终达成对多个目标域的通用适配。



（4）基于生成的方法

基于生成的域适应方法主要指的是利用图像到图像翻译（Image-to-Image Translation）的模型框架，将生成式模型直接作用于像素空间，
通过“以译代适”的思路缓解源域与目标域在颜色、纹理、噪声和光照等像素级统计上的差异，其核心通常由生成器和判别器（或共享编码器–多解码器结构）构成：
生成器负责将带标签的源域图像映射为目标域风格，同时尽可能保持语义内容与类别标签不变；判别器在图像或特征层面区分输入来自源域还是目标域，
通过对抗损失等推动生成器合成逼真且“目标域化”的图像，随后分类或分割等任务网络在这些“目标域风格源域样本”上进行监督训练，从而提升在真实目标域上的识别性能。
早期工作如PixelDA\cite{bousmalis2016unsupervised}和SimGAN\cite{shrivastava2017learning}，
利用生成对抗网络（Generative Adversarial Networks，GAN）\cite{goodfellow2014generative}对合成图像进行像素级转换，辅以内容保持或自正则项，在合成到真实数字、手势等任务上缩小了性能差距；
DTN\cite{taigman2016unsupervised}和UNIT\cite{liu2017unsupervised}进一步基于共享潜在空间假设学习双向无配对图像翻译，在此过程中获得可迁移的跨域内容表征。
而在存在显著结构差异或标签分布偏移时，单纯依赖图像到图像翻译易出现语义扭曲和类边界错配。

随后，研究焦点从“仅像素对齐”转向“像素+特征+语义任务”的联合自适应：CyCADA\cite{hoffman2018cycada}将CycleGAN\cite{zhu2017unpaired}式双向图像翻译与特征层对抗对齐、任务损失耦合，
系统消融表明多层次联合对齐效果最佳；Murez等人\cite{murez2017image}将无配对图像翻译嵌入共享编码器–解码器结构，要求同一编码器既能重建源域也能重建目标域图像，并结合特征域对抗约束学习域不变表示。
为更充分利用未标注目标域并缓解单向映射不稳定的问题,SBADA-GAN\cite{russo2017source}和DupGAN\cite{hu2018domain}通过双向或多域生成器同时建模源到目标和目标到源（或多源）的映射，引入类别一致性约束与域代码，
实现多域风格切换和潜在空间对齐；Li等人\cite{li2020generating}的目标域图像–标签对生成方法通过标签条件生成与语义损失显式对齐类条件分布，提升了生成图像的语义一致性和分类性能。
类似思想也被拓展到医学影像与遥感等领域：Benjdira等人\cite{benjdira2019unsupervised}和Li等人\cite{li2020evaluating}则利用CycleGAN等风格迁移模型在跨城市、
跨传感器遥感语义分割与目标检测中进行图像级域适应，系统评估表明在以光谱或外观差异为主的场景中像素级风格迁移有助于改善泛化性能。





\subsubsection{域泛化的研究现状}

域泛化任务可定义如下：源域数据集集合 $\mathcal{D}_s = \{D_{s_k} = \{(x_i^{s_k}, y_i^{s_k})\}_{i=1}^{n_{s_k}}\}_{k=1}^{K}$ 由 $K$ 个不同的源域构成，
其中每个源域 $D_{s_k}$ 均包含 $n_{s_k}$ 个带有标注的样本 $(x_i^{s_k}, y_i^{s_k})$；
目标域数据 $\mathcal{D}_t$ 代表未知的测试环境，在模型训练阶段完全不可见（Unseen）。
域泛化的核心任务是学习一个具有强泛化能力的映射函数 $f: \mathcal{X} \to \mathcal{Y}$，
仅利用 $K$ 个可见的源域数据集 $\mathcal{D}_s$ 进行训练，以实现对未知的目标域 $\mathcal{D}_t$ 
中样本的准确标签预测 $\hat{y} = f(x^t)$。需注意的是，
源域之间以及源域与未知目标域之间的概率特征分布均存在差异
（即 $P_{s_i}(X) \neq P_{s_j}(X)$ 且 $P_{s_k}(X) \neq P_t(X)$，其中 $i \neq j$）。
域泛化问题的设置如图\ref{fig:DG}所示，
基于深度学习的域泛化方法主要可划分为三类：基于数据增强的方法、基于表征学习的方法以及基于学习策略的方法。

\begin{figure}[htp]
\centering
\includegraphics[width=0.72\textwidth, keepaspectratio]{chap01/域泛化示意图.pdf}
\caption{域泛化问题建模概念示意图}
\label{fig:DG}
\end{figure}


（1）基于数据增强的方法

基于数据增强的域泛化方法的核心逻辑在于“模拟未知”：既然目标域不可见，
那么通过在训练过程中人为地增加源域数据的多样性，尽可能覆盖潜在的分布变化，理论上便能提升模型对未见域的鲁棒性。
这种增强通常发生在图像层面（Image-level）或特征层面（Feature-level）。
Tobin等人\cite{tobin2017domain}提出的域随机化（Domain Randomization）技术是此类方法的先驱，
通过在模拟环境中随机扰动纹理、光照和背景等参数生成海量合成数据，使得模型在无需真实数据的情况下也能具备良好的泛化能力。
在图像层面，除了传统的几何变换（如旋转、裁剪）和光度变换（如色彩抖动）以外，
基于生成式模型的方法占据了重要地位。Shankar等人\cite{shankar2018generalizing}利用贝叶斯数据增强技术，
通过扰动输入图像来生成具有挑战性的新样本，从而扩充源域的覆盖范围。随着生成对抗网络的发展，
Zhou等人\cite{zhou2020deep}提出了深度域对抗图像生成（Deep Domain-Adversarial Image Generation，DDAIG），
通过训练一个生成器来合成难以被当前分类器识别的新域数据，促使分类器学习更稳定的特征。此外，混合（Mixup）技术也被广泛引入DG领域，
Xu等人\cite{xu2020adversarial}提出的对抗性域混合技术（Adversarial Domain Mixup），通过在不同源域样本之间进行对抗性的插值混合，
生成位于域流形（Manifold）边缘的困难样本，从而扩展模型的泛化边界。
在特征层面，增强策略旨在直接干扰中间层表征以模拟风格变化。
Li等人\cite{li2021simple}发现不同域的特征统计量（均值和方差）通常携带了风格信息，
因此提出了特征统计量混合策略，通过随机交换不同样本的特征统计量来生成具有新颖风格特征的样本。
Zhou等人\cite{zhou2021domain}进一步提出了MixStyle方法，灵感源自风格迁移，通过在网络的浅层混合实例级特征统计量，模拟未见过的域风格，
且计算开销极低。 近年来，随着扩散模型（Diffusion Models）的兴起，基于扩散模型的增强方法也开始崭露头角。
Trabucco等人\cite{trabucco2023effective}展示了利用预训练的大规模扩散模型（如Stable Diffusion）进行以数据为中心的增强，
能够生成比传统GAN更丰富的跨域样本，并改善开放环境下的泛化性能。
总体而言，基于数据增强的方法通过扩展训练分布缓解了过拟合问题，
但如何保证增强数据的语义一致性以及如何平衡增强强度与训练稳定性，
仍是该方向需要进一步研究的问题\cite{wang2022generalizing}。



（2）基于表征学习的方法

基于表征学习的域泛化方法致力于寻找一种“最大公约数”式的特征表示：
即学习一个特征空间，使得在该空间内，所有源域的数据分布尽可能对齐，
同时保持类别可分性。其核心假设是：如果模型能够提取出源域间共享的域不变特征，那么这些特征在未知的目标域上同样适用。
该方向最经典的做法是显式地最小化不同源域分布之间的差异。
Muandet等人\cite{muandet2013domain}最早提出了域不变分量分析（Domain-Invariant Component Analysis，DICA），
通过核方法最小化分布差异。在深度学习时代，Li等人\cite{li2018domain}提出了对抗性特征学习框架，利用对抗训练的思想，
让特征提取器欺骗域判别器，从而学习到无法区分来源的通用特征。除了对抗学习，基于矩匹配（Moment Matching）的方法也很流行，
例如Peng等人\cite{peng2019moment}通过对齐特征的高阶矩来约束分布。
为了更精细地解耦特征，解耦表征学习（Disentangled Representation Learning）成为了近年来的研究热点。
该类方法试图将特征分解为“域特定特征”（Domain-Specific）和“域不变特征”（Domain-Invariant）。
Chattopadhyay等人\cite{chattopadhyay2020learning}提出了一种基于变分信息瓶颈（Variational Information Bottleneck）的框架，
旨在最大化预测准确率的同时，最小化特征中包含的特定域信息。此外，因果推断（Causal Inference）理论的引入为表征学习提供了全新的视角。
Arjovsky等人\cite{arjovsky2019invariant}提出的不变风险最小化（Invariant Risk Minimization，IRM），
试图寻找一种在所有环境下都能诱导最优分类器的特征表示，从因果关系而非统计相关性的角度定义泛化。
基于此，Mahajan等人\cite{mahajan2021domain}提出的MatchDG，通过匹配跨域样本的正负对来近似因果不变性，
虽然在复杂高维数据上的应用仍有挑战，但为解决虚假相关性（Spurious Correlation）提供了理论支撑。
这类方法是目前域泛化研究中最活跃的方向，但在面对源域数量较少或源域间差异过大的情况下，
强行对齐分布可能会导致负迁移或语义信息的丢失\cite{zhao2020domain}。


（3）基于学习策略的方法

基于学习策略的方法侧重于通过改进训练范式（如元学习、集成学习等）来增强模型的泛化潜能，
而非直接对数据或特征进行显式操作。这类方法通常模拟“训练-测试”的跨域过程，使模型“学会如何泛化”。
Li等人\cite{li2018learning}首次将元学习（Meta-Learning）引入域泛化，提出了MLDG（Meta-Learning for Domain Generalization）框架。
该方法将源域数据动态划分为元训练集（Meta-train）和元测试集（Meta-test），模拟真实的跨域场景：在元训练集上更新参数，
并要求更新后的模型在元测试集上也能表现良好，通过二阶梯度的优化直接引导模型寻找泛化极值点。 
Dou等人\cite{dou2019domain}进一步将此思想扩展，
提出了基于模型无关元学习（MAML）的MASF（Model-Agnostic Learning of Semantic Features），引入了全局类对齐正则化项，
在元学习过程中显式约束语义特征的一致性。另一种重要的策略是集成学习（Ensemble Learning），
Mancini等人\cite{mancini2018best}提出的预测模型通过学习每个源域特定的分类器，并在测试时根据输入样本与各源域的相似度加权融合预测结果，
以此利用特定域的专家知识来处理未知域。除此之外，自监督学习（Self-Supervised Learning）也被引入作为辅助任务，
Carlucci等人\cite{carlucci2019domain}通过引入一种名为“Jigsaw Puzzle”的自监督任务，
强迫模型在学习分类的同时理解图像的全局结构和空间关系，这种多任务学习策略被证明能有效防止模型过拟合于源域的表面纹理特征。
基于学习策略的方法\cite{balaji2018metareg,du2020learning}为域泛化提供了更加灵活的优化视角，
特别是元学习范式已成为提升模型在新环境中快速适应能力的标准工具之一\cite{wang2022generalizing}。



\subsection{跨域SAR目标识别的研究现状}
早期的SAR-ATR系统主要依赖于人工设计的特征提取与分类器构建。
这一阶段的研究核心在于如何从受斑点噪声干扰的 SAR 图像中提取鲁棒的物理与统计特征。
在特征提取层面，研究者利用 SAR 图像的后向散射特性，开发了多种特征描述子。
几何特征（如目标轮廓、阴影形状）被广泛用于描述目标的物理结构 \cite{sun2007adaptive,papson2012classification}。
基于幅度的统计特征（如灰度共生矩阵、K分布参数）则侧重于捕捉目标的纹理信息 \cite{fukuda2001support}。
此外，为了利用 SAR 信号的电磁散射机理，基于属性散射中心（Attributed Scattering Center, ASC）的模型被提出，
通过提取目标的散射点拓扑结构来实现精确描述 \cite{duan2014polarimetric,he2014forward}。变换域特征，如小波变换和单演信号特征，
也被用于增强目标在不同尺度下的表达能力 \cite{achim2003sar}。
在分类器设计层面，传统的模式识别算法占据主导地位。
最近邻分类器和支持向量机因其理论成熟而被广泛应用\cite{zhao2002support,blanzieri2008nearest}。
其中，基于稀疏表示的分类方法通过在过完备字典上重构测试样本，
有效提升了噪声环境下的识别鲁棒性 \cite{dong2015sar}。尽管传统方法在诸如MSTAR数据集的SOC这类简单场景下表现尚可，
但在扩展操作条件（Extended Operating Conditon，EOC）下（如俯仰角变化、背景变化），由于人工特征的表达能力有限，其泛化性能往往面临瓶颈。

随着深度学习技术的兴起，基于卷积神经网络（Convolutional Neural Network，CNN）的方法逐渐成为 SAR-ATR 的主流。
深度模型能够从原始数据中自动学习分层特征，大幅提升了识别精度与泛化能力。
经典 CNN 模型的迁移与改良是研究的起步点。
研究者将 AlexNet、ResNet 等在光学图像上成功的模型迁移至 SAR 领域，
并针对 SAR 数据量小的特点进行轻量化改进。例如，Chen 等人提出的全卷积网络（A-ConvNets）
通过移除全连接层显著减少了参数量，有效缓解了过拟合问题\cite{chen2016target}。
为了进一步提升复杂场景下的识别性能，结合领域知识的深度学习成为热点。
部分研究尝试将传统物理特征（如 ASC 模型）嵌入神经网络，构建物理可解释的深度网络\cite{feng2021sar,zhang2025masff}。
注意力机制（Attention）的引入，使得网络能够聚焦于目标的关键散射区域，抑制背景杂波的干扰\cite{zhang2020convolutional}。
此外，针对多视角、多模态数据的特征融合策略也被证明能有效提高识别的置信度\cite{pei2017sar}。
当前，跨域识别是该领域的前沿挑战，由于训练与测试数据往往来自不同传感器或成像条件，域偏移问题显著影响模型性能。
为此，研究者们引入跨域迁移学习技术，以提升模型的跨域鲁棒性。

\subsubsection{域适应SAR目标识别}
迁移学习与SAR图像目标识别的结合应用起步较早，其核心目标是解决SAR数据样本稀缺场景下深度模型易发生过拟合的关键问题。
鉴于SAR数据获取成本较高、样本规模相对有限的固有特性，
以及深度神经网络浅层特征具备的强通用性与跨域迁移潜力，
微调（fine-tune）技术成为该场景下的经典解决方案。
具体而言，该方法通常先在大规模光学数据集或分布相关的SAR辅助数据集上完成模型预训练，
利用大量数据学习通用视觉特征表示；随后固定预训练模型的浅层网络参数（避免通用特征被破坏），
仅采用目标 SAR 识别任务的少量标注数据对网络深层结构进行微调优化，
通过适配特定任务的特征空间，实现模型识别性能的有效提升。
在相关研究实践中，Kang等人\cite{kang2016sar}于光学数据集CIFAR-10上完成深度神经网络预训练后，
固定网络浅层参数并仅对最深两层网络进行微调，该策略显著提升了模型对TerraSAR图像的场景分类精度；
Malmgren等人\cite{malmgren2017improving}则采用电磁散射仿真生成的SAR车辆目标图像构建预训练数据集，
基于该数据完成模型初始化后，利用实测车辆目标图像对网络进行微调优化，最终实现了目标识别性能的显著改善，
也验证了仿真数据预训练-实测数据微调范式的有效性。

然而，预训练-微调范式并非在所有跨域场景中均能有效提升SAR目标识别性能。
Girshick等人\cite{girshick2014rich}的研究证实，即使经过预训练与微调处理，
源域与目标域之间的分布偏移仍会对深度网络的目标识别性能产生显著负面影响。
这一现象的核心原因在于，深度神经网络的特征表示具有层级演进特性：浅层特征以低级视觉信息（如边缘、纹理）为主，
迁移性较强；而深层特征逐渐趋向抽象化与任务特异性，其跨域迁移能力随网络层级加深而急剧下降。
因此，直接将在源域数据集上学习的深层特征迁移至目标域数据集时，往往会因分布不匹配导致分类性能大幅退化。
针对这一问题，当目标域数据仅部分可见时，众多研究者开始引入域适应方法，
为解决跨域SAR目标识别中的分布偏移问题提供了新的技术路径。

Huang等人\cite{huang2019and}深入探究了自然图像数据集ImageNet与不同来源SAR图像间深层特征的迁移性能，
发现直接采用ImageNet数据预训练OpenSARShip舰船分类模型会引发负迁移问题。
为此，他们提出一种基于多核极大均值差异（Multikernel Maximum Mean Discrepancy，MK-MMD）的域自适应方法，通过匹配多源数据间的特征分布，有效提升了模型的预训练效率。
针对不同成像条件下的SAR图像目标识别任务，Shi等人\cite{shi2021unsupervised}为SAR目标检测任务构建了融合生成对抗与对抗学习的域自适应网络，
分别对浅层、深层及样本层特征实施域适配，并通过选择性构建目标域伪标签的方式实现网络的迭代训练。
Chen等人\cite{chen2022pixel}提出一种两阶段无监督域自适应方法，该方法融合风格迁移网络与对抗学习机制：
首先借助CycleGAN将源域SAR图像转换为目标域图像风格，实现像素级的对齐，
随后通过生成图像与真实图像间的对抗学习，完成特征分布的对齐。

部分研究者将仿真SAR图像引入SAR图像目标识别任务，以缓解实测数据样本稀缺的问题。
He等人\cite{he2021sar}针对训练与测试数据成像条件差异导致深度学习模型识别性能下降的问题，
提出一种基于模拟数据的任务驱动域适应方法，通过引入仿真SAR数据并结合MK-MMD域混淆度量最小化域差异，
在MSTAR数据集的不同扩展操作条件下均取得了优于其他对比方法的识别性能。
针对仿真与实测图像存在的视觉差异，Youk等人\cite{youk2023transformer}设计了一个基于Transformer的仿真SAR图像到实测SAR图像风格迁移模型。
其模型仅需部分类别的成对仿真和实测图像进行训练，在未见过的类别上仍有良好的迁移能力，
最终甚至可以实现在部分类别上利用纯仿真数据训练实测SAR目标识别模型。
Shi等\cite{shi2024unsupervised}提出一种基于域级和类级双粒度对齐的无监督域自适应方法，通过梯度加权对抗学习实现域级全局对齐、
原型网络结合对比学习建模跨域原型关系实现类级对齐、属性散射中心辅助伪标签过滤提升可靠性，
显著提升了跨仿真-实测域SAR目标识别性能。

综上，现有的研究大多聚焦于无监督域适应场景，有效缓解了域偏移问题，但模型训练过程仍需访问相当数量的目标域实测SAR图像。
本质上解决的还只是SAR图像标注困难的问题，无法解决实际应用中目标域SAR数据难以获取的挑战。
在现有域适应研究基础上，探索目标域数据少样本条件下的域适应SAR目标识别方法具有重要的理论意义和应用价值。



\subsubsection{域泛化SAR目标识别}
长期以来，由于SAR图像易变的特性，如何提升SAR目标识别模型的鲁棒性一直都是研究的热点问题，也取得了可观的研究成果。
例如在解决MSTAR数据集的扩展工作条件方面，
Li等人\cite{li2023hierarchical}针对SAR车辆识别中存在类内变异大、类间变异小、数据集缺乏的问题，
且现有方法易受背景杂波干扰、域对齐未充分考虑目标与杂波分离的不足，提出一种分层解纠缠-对齐网络，
通过域数据生成模块、多任务辅助掩码解纠缠模块、目标特征域对齐模块的协同作用，实现不同操作条件下的稳健识别。
Xiong等人\cite{xiong2024lightweight}针对SAR-ATR中训练数据稀缺易过拟合、模型存在黑箱性、
计算需求高且无人机平台对算法轻量化和适应性有更高要求的问题，提出一种轻量级双流复数SAR目标识别框架，
该框架由基于多级多头注意力机制引导的异构图网络，通过ASC模型提取电磁散射信息构建异构图，
学习局部电磁散射与拓扑特征和基于SE-ResNet18的全局视觉特征提取网络组成，
借助特征融合子网自适应融合双流特征，并结合分类损失与拓扑损失提升类内特征聚合性，兼顾识别鲁棒性与轻量化需求。
随着计算机视觉基础大模型的发展，Li等人\cite{li2025saratr}针对SAR-ATR模型存在任务特定性强、过度依赖昂贵标注、
模型未适配SAR图像特性、缺乏大规模开源数据集的核心问题，提出首个SAR-ATR基础模型 SARATR-X。
该模型基于14个开源数据集构建含18.66万样本的预训练数据集SARDet-180K，采用适配小目标的 HiViT 骨干网络，
通过 “ImageNet 预训练初始化+SAR图像多尺度梯度特征抑制斑点噪声” 的两步自监督策略学习通用表征，
并用于少样本、稳健性分类及多类别、特定类别检测任务。

专注提升模型鲁棒性虽然有效提升了SAR目标识别模型的综合性能，
但一方面在训练数据与测试数据存在更显著分布差异的场景中模型泛化能力仍有所不足——比如以纯仿真SAR图像训练，
而在实测图像测试的场景，另一方面对模型结构的复杂设计也增加了实际部署应用的难度。
为切实稳定地提升模型在未知数据分布上的识别性能，少部分研究人员将研究重心转向域泛化方法，
主要研究域泛化训练策略以实现跨域识别，避免对骨干网络的大幅度改动。
Inkawhich\cite{inkawhich2021bridging}等提出了在仅使用仿真SAR数据训练时提高分类模型泛化性能的技术，其工作表明
模型在实测数据上的分类性能会根据骨干网络的类型、数据增强技术（如高斯噪声、混合和旋转）、损失函数及其组合而变化。
其通过实验和分析筛选组合策略，并通过集成学习将五种仅通过仿真数据训练的分类模型集成，
进而在实测样本上获得较好的分类结果。
Kim等\cite{kim2024soft}提出一种软分割域随机化方法，利用高斯混合模型对合成 SAR 数据中的目标与杂波区域进行软分割，
并有选择性地对杂波进行随机化处理，从而缩小合成数据与真实数据的统计特性差异。
Jang等\cite{jang2025irasnet}提出了一种基于特征解耦的域泛化框架，能够有效实现特征级别的杂波减少和域不变特征学习。
其首先提出了一个杂波减少模块，该模块能在特征图上最大化信号与杂波比率，在特征级别上减少杂波的影响，
同时保留了目标和阴影信息。其次，其将对抗性学习与杂波减少模块整合，以提取减少杂波后的域不变特征。
Yuan等人\cite{yuan2024filling}提出了一个单源域泛化SAR舰船检测框架。
针对单一源域特征分布空间的限制，其提出了一种基于特征归一化的简单有效的插值方法，允许框架模拟不同设备捕获的图像，
使特征提取器能够学习到泛化的特征表示。另一方面，为了更好地训练一个更跨域稳定的检测器，
提出了一个基于对比学习的高阶实例解耦模块，通过实例级对比损失和熵损失作为额外的约束，
将原始特征分解为任务相关和任务无关特征。在两个层面上增强了SAR舰船检测模型的泛化能力。

综上所述，相较于自然图像领域，跨域SAR目标识别研究仍处于持续发展阶段。
围绕少样本域适应与域泛化等问题，现有工作已给出初步探索，但在数据受限条件下的性能稳定性、方法复杂度和实际部署可行性等方面仍有改进空间。
针对复杂多变的应用场景，进一步引入并改进域适应与域泛化技术，
有助于缓解实测SAR训练数据受限带来的性能瓶颈，
并提升数据驱动型SAR目标识别算法在实际场景中的稳定性与可用性。


\section{本文结构安排}

本文共分为五章，各章内容安排如下：

\textbf{第一章：绪论。}
本章首先阐述SAR目标识别的研究背景与意义，
分析实测训练数据获取困难、跨域数据分布偏移显著以及边缘平台部署资源受限等问题；
随后从跨域迁移学习技术和跨域SAR目标识别两个层面梳理国内外研究现状；
在此基础上明确本文所关注的“目标域数据少样本”和“目标域数据不可见”两类问题场景，以及相应的研究思路与主要研究内容；
最后给出全文的整体结构安排，为后续章节展开奠定基础。

\textbf{第二章：面向目标域数据少样本的域适应SAR目标识别方法。}
本章首先在引言中说明少样本目标域条件下跨域SAR目标识别问题的研究背景与本章思路；
随后围绕所采用的域适应方法展开介绍，先给出以像素级域对齐和类混淆域适应为核心的总体框架与主要流程，
再按照方法概述、像素级域对齐模块和类混淆域适应模块的顺序，对块噪声对比估计损失、随机数据增强一致性正则化以及目标域判别边界优化等内容进行阐述，
说明各组成部分的设计动机、输入输出关系以及训练过程；
在实验部分，本章首先介绍SAMPLE数据集及跨域实验任务设置，
然后说明评价指标、训练参数和对比方法等实验条件，
进一步通过对比实验、消融实验以及图像翻译结果和特征分布的可视化分析，对方法性能、模块作用和少样本条件下的域适应效果进行系统验证与讨论；
最后对本章内容进行总结。

\textbf{第三章：面向目标域数据不可见的域泛化SAR目标识别方法。}
本章首先在引言中说明目标域数据不可见条件下跨域SAR目标识别问题的研究背景与分析思路；
随后对域泛化方法的整体框架进行介绍，并按照方法概述、像素级数据增强模块、特征级风格混淆模块和域对抗学习模块的顺序展开论述，
依次说明高斯噪声注入、高频频域扰动、分贝化变换、基于实例归一化的风格混淆机制以及基于梯度反转层的域对抗学习等具体内容在整体方法中的作用、相互关系及实现方式；
在实验部分，本章先介绍SAMPLE和S2M-5数据集及跨域测试任务的设置，
再对评价指标、训练参数和实验流程进行说明，
并通过对比实验、消融实验以及数据增强效果与特征分布的可视化分析，对方法在不同数据集和不同设置下的表现进行验证，
从而分析其在目标域不可见条件下的适用性与稳定性；
最后对本章内容进行总结。

\textbf{第四章：面向国产AI芯片的跨域SAR目标识别模型轻量化部署。}
本章首先在引言中说明跨域SAR目标识别模型面向国产AI芯片平台部署的需求与研究内容；
随后从模型轻量化处理和端侧部署实现两个部分展开。
其中，模型轻量化部分主要介绍以第二章和第三章训练得到的ResNet50模型为教师网络、以MobileNetV2为学生网络的知识蒸馏过程，以及基于AMCT的INT8后训练均匀量化方案，
包括网络选型依据、蒸馏损失设计、训练流程与量化方案设计等；
部署实现部分主要介绍华为Atlas 200I DK A2平台环境、模型导出、AMCT量化、ATC转换、基于AscendCL的推理实现以及端侧推理链路优化，
对模型从训练结果到板端执行的完整部署过程进行说明。
在实验与分析部分，本章进一步通过蒸馏训练实验设置、轻量化前后识别性能对比、模型压缩与端侧效率分析以及综合讨论，
对两类跨域模型在轻量化和部署前后的精度保持情况、模型压缩效果和端侧运行性能进行验证；
最后对本章内容进行总结。

\textbf{第五章：总结与展望。}
本章对全文的研究工作进行系统总结，
归纳本文的主要研究内容与结论，
并结合跨域SAR目标识别与国产平台部署的发展需求，
从后续方法研究与工程应用两个层面对未来工作进行展望。

本文的总体结构安排如图\ref{fig:chap1_structure}所示。

\begin{figure}[htbp]
	\centering
	\includegraphics[width=\textwidth, keepaspectratio]{论文结构图v3.png}
	\caption{本文结构安排}
	\label{fig:chap1_structure}
\end{figure}


